\documentclass[../DoAn.tex]{subfiles}
\begin{document}

\begin{center}
    \Large{\textbf{ABSTRACT}}\\
\end{center}
\vspace{1cm}

Network intrusion detection requires classifying network traffic into benign and attack categories in real time. Signature-based systems such as Snort effectively detect known attacks but fail against previously unseen variants and behavioral exploits. Machine learning approaches can learn behavioral features from labeled data, yet face two practical challenges: severe class imbalance (benign traffic typically exceeds 80\% of all flows) and covariate shift when a model trained on a benchmark dataset is deployed in a different network environment.

This thesis develops a hybrid network intrusion detection system combining Snort and machine learning, built through three sequential research phases. Phase 1 establishes a foundational two-stage architecture on NSL-KDD: an Autoencoder Anomaly Gate filters normal traffic at the first stage, while a Stacking Ensemble of FT-Transformer and LightGBM combined via a Meta Logistic Regression classifies attack types at the second stage. The model achieves Macro F1 = 0.6809 on KDDTest+ under a strict train/test separation protocol. Phase 2 redesigns the architecture as a Two-Stage Cascade for CIC-IDS-2017: a Gating Network performs 6-class classification (Benign, DoS, DDoS, PortScan, Brute Force, Suspicious) at an 85\% Suspicious-confidence threshold, while an Expert Network classifies 5 output classes (Benign, Web Attack, Bot, Infiltration, Heartbleed) for flows forwarded from Stage 1; together they produce 9 distinct output labels. Two rounds of Hard Negative Mining with 10$\times$ sample replication outperform class-weight adjustment on the Botnet class (F1 = 0.6512 versus 0.5103). Asymmetric Ensemble Voting then combines FT-Transformer, Random Forest (150 trees, class\_weight=balanced), and KNN (K=16, uniform weighting) through a veto mechanism rather than majority voting, raising end-to-end Botnet F1 from 0.6294 to 0.7344 by improving its precision from 47.87\% to 79.47\%. The system achieves 99.55\% Accuracy and Macro F1 = 0.9294 on the CIC-IDS-2017 test set. Phase 3 addresses covariate shift when deploying the CIC model to a WSL2 testbed environment, where accuracy drops from 99.55\% to 21.93\% under direct transfer. After collecting real attack traffic from two machines on a home LAN, the final model FT-IDS is obtained by fine-tuning the Phase 2 backbone --- extended from 77 to 80 features through Model Surgery --- rather than training from scratch, and achieves Macro F1 = 91.66\% across 5 classes, with BruteForce Recall = 90\% on a domain-diverse validation set.

The main contributions are: (1) a Two-Stage Cascade architecture with Asymmetric Ensemble Voting for severely imbalanced multi-class intrusion detection; (2) evidence that two-round Hard Negative Mining outperforms class-weight adjustment on minority classes by concentrating the learning signal near the decision boundary; (3) a quantified covariate shift diagnosis together with a transfer strategy based on selective layer freezing and Model Surgery; and (4) empirical evidence that domain-homogeneous validation sets can inflate reported metrics by up to 14 percentage points compared to domain-diverse evaluation.

\end{document}
